% Μεμονωμένο Παράρτημα 
% ορίζεται ο τίτλος του παραρτήματος
\edef\chaptertitle{Τίτλος Παραρτήματος}
\label{appendix}

% Μην τροποποιήσετε τις ακόλουθες εντολές
\chapter*{\chaptername\\\vspace{1cm} \chaptertitle} 
\addcontentsline{toc}{chapter}{\chaptername\texorpdfstring{\\}~\chaptertitle}
\setcounter{section}{0}
\renewcommand{\chaptermark}[1]{\markboth{\textit{\chaptername. \chaptertitle}}{}}
\chaptermark{}


\section{Tables}

\begin{table}[h!]
	\centering
	\caption{Complete GPIO Pinout for Raspberry Pi 4 Model B (Part 1/2)}
	\label{table:rpi4_gpio_part1}
	\begin{tabular}{|c|l|l|c|l|}
		\hline
		\textbf{Pin \#} & \textbf{Function} & \textbf{Alt Function} & \textbf{Voltage} & \textbf{Notes} \\ \hline
		1 & 3.3V & -- & 3.3V & 50mA max \\ \hline
		2 & 5V & -- & 5V & Power supply \\ \hline
		3 & GPIO2 & SDA1 & 3.3V & I²C Data \\ \hline
		4 & 5V & -- & 5V & Power supply \\ \hline
		5 & GPIO3 & SCL1 & 3.3V & I²C Clock \\ \hline
		6 & GND & -- & -- & Ground \\ \hline
		7 & GPIO4 & GPCLK0 & 3.3V & General purpose \\ \hline
		8 & GPIO14 & TXD0 & 3.3V & UART Transmit \\ \hline
		9 & GND & -- & -- & Ground \\ \hline
		10 & GPIO15 & RXD0 & 3.3V & UART Receive \\ \hline
		11 & GPIO17 & -- & 3.3V & General purpose \\ \hline
		12 & GPIO18 & PCM\_CLK & 3.3V & PWM0 \\ \hline
		13 & GPIO27 & -- & 3.3V & General purpose \\ \hline
		14 & GND & -- & -- & Ground \\ \hline
		15 & GPIO22 & -- & 3.3V & General purpose \\ \hline

	\end{tabular}
\end{table}

\begin{table}[h!]
	\centering
	\caption{Complete GPIO Pinout for Raspberry Pi 4 Model B (Part 2/2) (Continued from Table \ref{table:rpi4_gpio_part1})}
	\label{table:rpi4_gpio_part2}
	\begin{tabular}{|c|l|l|c|l|}
		\hline
		\textbf{Pin \#} & \textbf{Function} & \textbf{Alt Function} & \textbf{Voltage} & \textbf{Notes} \\ \hline
		16 & GPIO23 & -- & 3.3V & General purpose \\ \hline
		17 & 3.3V & -- & 3.3V & 50mA max \\ \hline
		18 & GPIO24 & -- & 3.3V & General purpose \\ \hline
		19 & GPIO10 & MOSI0 & 3.3V & SPI0 \\ \hline
		20 & GND & -- & -- & Ground \\ \hline
		21 & GPIO9 & MISO0 & 3.3V & SPI0 \\ \hline
		22 & GPIO25 & - & 3.3V & General purpose \\ \hline
		23 & GPIO11 & SCLK0 & 3.3V & SPI0 \\ \hline
		24 & GPIO8 & CE0 & 3.3V & SPI0 Chip Enable 0 \\ \hline
		25 & GND & - & - & - \\ \hline
		26 & GPIO7 & CE1 & 3.3V & SPI0 Chip Enable 1 \\ \hline
		27 & ID\_SD & - & - & EEPROM ID \\ \hline
		28 & ID\_SC & - & - & EEPROM Clock \\ \hline
		29 & GPIO5 & - & 3.3V & General purpose \\ \hline
		30 & GND & - & - & - \\ \hline
		31 & GPIO6 & - & 3.3V & General purpose \\ \hline
		32 & GPIO12 & PWM0 & 3.3V & PWM channel 0 \\ \hline
		33 & GPIO13 & PWM1 & 3.3V & PWM channel 1 \\ \hline
		34 & GND & - & - & - \\ \hline
		35 & GPIO19 & PCM\_FS & 3.3V & PWM/MISO1 \\ \hline
		36 & GPIO16 & - & 3.3V & General purpose \\ \hline
		37 & GPIO26 & - & 3.3V & General purpose \\ \hline
		38 & GPIO20 & MOSI1 & 3.3V & SPI1 \\ \hline
		39 & GND & - & - & - \\ \hline
		40 & GPIO21 & MISO1 & 3.3V & SPI1 \\ \hline
\end{tabular}
\end{table}

\clearpage 






\subsection{Αλγόριθμοι}
Ακολουθεί ο Αλγόριθμος~\ref{alg1}, ο οποίος είναι μορφοποιημένος με τα πακέτα algorithm και algorithmic.

\begin{algorithm}[htb]
\caption{\ \ \ Probabilistic $k\theta NN$ Monitoring}
\begin{algorithmic}[1]
\begin{small}

\STATE{\bf Procedure} {\em VerifyCandidate} (focal query point $q$, threshold $\theta$, object $o$, list of auxiliary objects $P$, distance $kMAXDIST$) 

\IF { $\Phi(o, kMAXDIST) \geq \theta$ {\bf and} $L_2(q, o) \leq L_2(q, P.$top()) }

\STATE {$P$.pop()};   \ \ \ \ \ \ \ \ {\em //Replace the most extreme element in $P$, since candidate $o$ ... }

\STATE {$P$.push($o$)};  \ \ \ \ \ {\em //... has enough probability and has its mean closer to focal $q$ }

\ENDIF

\STATE {\bf End Procedure}


\end{small}
\end{algorithmic}
\label{alg1}
\end{algorithm}


\subsection{Θεωρήματα, Πορίσματα, Ορισμοί, κλπ.}

Ακολουθεί παράδειγμα θεωρήματος από την ιστοσελίδα  {\small\url{https://www.overleaf.com/learn/latex/Theorems_and_proofs}}

\begin{theorem}
Let $f$ be a function whose derivative exists in every point, then $f$ 
is a continuous function.
\end{theorem}



